% Draft methods subsection for the dynamic evaluation.
% Every number here is reproducible from the repository:
%   python -m sphyr.evaluate_dynamic validate --subject full_easy --samples 20
%   python -m sphyr.evaluate_dynamic rescore --workers 8

\subsection{Dynamic Evaluation}
\label{sec:dynamic-evaluation}

A topology optimisation problem does not have a unique solution. Comparing a
completion to the reference answer cell by cell therefore measures the wrong
thing twice over: a structure that routes the load differently but just as
efficiently is marked wrong, while a structure that matches the reference almost
everywhere but severs the load path is marked nearly right. Across the
\num{24325} valid completions in our results, \SI{72.7}{\percent} carry their
load correctly while differing from the reference answer somewhere.

We therefore evaluate completions \emph{dynamically}: each one is simulated as a
structure, and scored against a topology optimiser re-run on the same masked
region under the same material budget.

\paragraph{Structural model.}
Each grid cell is a unit square bilinear (Q4) plane-stress finite element with
Young's modulus interpolated from the cell density by the SIMP law
\begin{equation}
  E(\rho) = E_{\min} + \rho^{\,p}\,(E_0 - E_{\min}),
  \qquad p = 3,\; E_0 = 1,\; E_{\min} = 10^{-9},
  \label{eq:simp}
\end{equation}
so a density of $0$ leaves a numerically negligible but non-singular stiffness
and a density of $1$ recovers the full material stiffness. The load cells
($\mathtt{L}$) carry a unit resultant spread evenly over their nodes, and the
support cells ($\mathtt{S}$) are clamped in both degrees of freedom; load and
support cells are solid by definition. Boundary conditions are always read from
the reference sample, so a completion cannot improve its score by relocating
loads or supports. Solving $\mathbf{K}\mathbf{u} = \mathbf{f}$ gives the
\emph{compliance} $c = \mathbf{f}^{\!\top}\mathbf{u}$, the work done by the load
and hence the inverse of stiffness.

Compliance also decides, physically rather than heuristically, whether a
structure carries its load at all: a completion whose compliance exceeds
$10^{3}$ times that of the fully solid grid is transmitting force through the
void stiffness rather than through material. This is strictly sharper than a
connectivity test on the cell graph. On our results the two disagree on
\num{1079} completions: \num{486} that graph connectivity calls connected but
that cannot transmit the load---corner-touching cells share a single node and
form a hinge---and \num{593} that carry load despite the graph calling them
disconnected.

\paragraph{Reference design.}
For each sample we re-run a SIMP optimiser~\cite{bendsoe2003topology} on
\emph{exactly} the masked cells, with every unmasked cell pinned to the value
the sample prescribes, minimising compliance subject to the reference material
budget. The optimiser uses a cone density filter of radius $1.2$, an
optimality-criteria update, penalisation continuation over $p \in \{1,2,3\}$,
volume-preserving binarisation, and a discrete swap refinement. Because a single
SIMP run on a $10\times10$ mesh has grey local minima, we generate several
candidate designs---cold and warm-started runs at each penalisation, their
binarisations, and the reference answer itself---and keep the stiffest one that
respects the budget. Seeding the search with the reference answer guarantees
that the design a completion is measured against is never worse than the answer
the dataset ships, which we verify on every run.

\paragraph{Score.}
Let $c$ and $v$ be the compliance and volume fraction of a completion, $c^\star$
the compliance of the reference design, and $v_{\mathrm{ref}}$ the reference
material budget. We report
\begin{align}
  \mathrm{SE} &= \min\!\left(1, \; c^\star / c\right) & \text{(structural efficiency)} \label{eq:se}\\
  \mathrm{ME} &= \min\!\left(1, \; v_{\mathrm{ref}} / v\right) & \text{(material efficiency)} \label{eq:me}\\
  \mathrm{TS} &= \mathrm{SE}\cdot\mathrm{ME} & \text{(topology score)} \label{eq:ts}
\end{align}
with $\mathrm{SE} = \mathrm{ME} = \mathrm{TS} = 0$ for a completion that carries
no load.

Both factors are necessary. Stiffness alone is maximised by filling the grid
solid, and thrift alone by building nothing; only the product rewards carrying
the load with the material the task allows. The need for the second factor is
also visible empirically. Among load-carrying completions, \SI{68.5}{\percent}
of those that spend more material than the reference achieve a lower compliance
than the budget-constrained reference design---as they should, since they are
spending more. Restricted to completions that spend the same material, that
figure falls to \SI{5.3}{\percent}, and to \SI{0.3}{\percent} for those spending
less. Equation~\eqref{eq:se} on its own would read material overspend as design
quality.

\paragraph{Validation.}
We check that the score orders completions whose relative quality is known in
advance. On \texttt{full\_easy} (20 samples), the reference answer scores
$0.93$, the same answer with two cells flipped $0.80$, with ten cells flipped
$0.48$, a grid filled solid $0.29$, and an empty grid $0.00$. The reference
answer does not score $1.0$ because the optimiser usually finds a slightly
stiffer structure for the same material; that margin is the headroom a model is
measured against. Across every validation run the reference design was never
beaten by the reference answer it bounds.
